
\subsection{משוואת לפלס בגיאומטריות שונות}
\textbf{תאריך ההרצאה: 13.01.2026}

\subsection{מבוא פיזיקלי ומתמטי}
משוואת לפלס:
\[
\nabla^2 u = 0
\]
מתארת פונקציות \emph{הרמוניות} (\LR{Harmonic Functions}), המופיעות במגוון הקשרים פיזיקליים:
פוטנציאל חשמלי בהיעדר מטענים, הטמפרטורה הסטציונרית במשוואת הדיפוזיה, זרימה פוטנציאלית ועוד.

\begin{mainbox}{עקרונות שיטת הפתרון}
העיקרון המרכזי המאפשר את הפתרון הוא הליניאריות:
\begin{itemize}
    \item המשוואה \textbf{ליניארית והומוגנית}.
    \item סכום (ואף טור) של פתרונות הוא פתרון (סופרפוזיציה).
    \item פתרון באמצעות \textbf{הפרדת משתנים} נותן משפחה של פתרונות פרטיים (פונקציות עצמיות), שמהם נבנה הפתרון הכללי.
\end{itemize}
\end{mainbox}

\subsection{משוואת לפלס בקואורדינטות פולריות}
נניח סימטריה גלילית ארוכה (אין תלות ב-\(z\)). אזי הלפלסיאן הוא:
\[
\nabla^2 u =
\frac{1}{r}\frac{\partial}{\partial r}
\left(r\frac{\partial u}{\partial r}\right)
+ \frac{1}{r^2}\frac{\partial^2 u}{\partial \theta^2}.
\]

נבצע הפרדת משתנים \(u(r,\theta)=R(r)\Theta(\theta)\) ונקבל:
\[
\frac{r^2 R'' + r R'}{R} = -\frac{\Theta''}{\Theta} = \lambda^2.
\]

\subsection{המשוואה הזוויתית והערכים העצמיים}
המשוואה הזוויתית היא:
\[
\Theta'' + \lambda^2 \Theta = 0.
\]

בבעיות עם מחזוריות מלאה בזווית (עיגול מלא, טבעת) נדרש תנאי שפה מחזורי \(\Theta(\theta+2\pi)=\Theta(\theta)\). לכן הערכים העצמיים הם שלמים:
\[
\lambda = n,\qquad n\in\mathbb{N}_0.
\]

הפתרונות:
\[
\Theta_n(\theta)=
\begin{cases}
1, & n=0,\\
\cos(n\theta),\ \sin(n\theta), & n\ge 1.
\end{cases}
\]

\subsection{המשוואה הרדיאלית}
ל-\(n\ge1\) מתקבלת משוואת אוילר (\LR{Euler Equation}):
\[
r^2 R'' + r R' - n^2 R = 0,
\]
עם פתרונות:
\[
R_n(r)=C_1 r^n + C_2 r^{-n}.
\]

ל-\(n=0\):
\[
r^2 R'' + r R' = 0
\quad\Rightarrow\quad
R_0(r)=C_1 \ln r + C_2.
\]

\subsection{בעיה פנימית במעגל}
תחום: \(0\le r\le R\).

\begin{notebox}{חסימות בראשית}
בבעיה פנימית הכוללת את הראשית, הפתרון חייב להיות \textbf{חסום ב-\(r=0\)}.
לכן כל איבר שמתבדר ב-\(r\to0\) נפסל (כלומר \(r^{-n}\) ו-\(\ln r\)).
\end{notebox}

מכאן:
\[
R_n(r)=r^n,\quad n\ge1,
\qquad
R_0(r)=\text{קבוע}.
\]

\begin{resultbox}{הפתרון הכללי בעיגול}
\[
u(r,\theta)=
A_0
+\sum_{n=1}^\infty
r^n\bigl(
A_n\cos(n\theta)+B_n\sin(n\theta)
\bigr).
\]
\end{resultbox}

\subsection{קביעת המקדמים: טור פורייה}
נתון תנאי שפה דיריכלה על השפה: \(u(R,\theta)=G(\theta)\).
מאחר ש-\(G\) מחזורית ב-\(2\pi\), נפרוס אותה בטור פורייה:
\[
G(\theta)=
\frac{a_0}{2}
+\sum_{n=1}^\infty
\bigl(
a_n\cos(n\theta)+b_n\sin(n\theta)
\bigr).
\]

השוואת מקדמים נותנת:
\[
A_n=\frac{a_n}{R^n},
\qquad
B_n=\frac{b_n}{R^n},
\qquad
A_0=\frac{a_0}{2}.
\]

\subsection{דוגמה}
נתון תנאי השפה:
\[
u(R,\theta)=
3\sin\theta -2\cos(3\theta).
\]

מייד מתקבל הפתרון (ללא צורך באינטגרלים):
\[
u(r,\theta)=
3\frac{r}{R}\sin\theta
-2\frac{r^3}{R^3}\cos(3\theta).
\]

\begin{examplebox}{אורתוגונליות}
אין צורך לחשב מקדמי טור פורייה באמצעות אינטגרלים כאשר תנאי השפה כבר נתון כסכום מפורש של פונקציות טריגונומטריות (אופני לפלס).
\end{examplebox}

\subsection{בעיה חיצונית למעגל}
תחום: \(r\ge R\).

\begin{notebox}{חסימות באינסוף}
בבעיה חיצונית נדרש בדרך כלל שהפתרון יהיה \textbf{חסום באינסוף} (\(r\to\infty\)),
אלא אם נאמר במפורש אחרת (מקורות באינסוף).
\end{notebox}

לכן עבור \(n\ge1\) נבחר \(R_n(r)=r^{-n}\). עבור \(n=0\), אם דורשים חסימות, נשאר קבוע בלבד.

הפתרון הכללי:
\[
u(r,\theta)=
A_0
+\sum_{n=1}^\infty
r^{-n}\bigl(
A_n\cos(n\theta)+B_n\sin(n\theta)
\bigr).
\]

\begin{physicsnote}{משמעות האיבר הלוגריתמי}
אם \emph{לא} נדרש שהפוטנציאל יהיה חסום באינסוף, האיבר \(C\ln r\) מותר בפתרון ל-\(n=0\). פיזיקלית, הוא מייצג פוטנציאל של תיל טעון אינסופי (\LR{Line Charge}). בבעיות מעבר חום (טמפרטורה סטציונרית), איבר זה בדרך כלל נפסל משיקולי אנרגיה או טמפרטורה סופית.
\end{physicsnote}

\subsection{בעיית הטבעת (Annulus)}
תחום: \(R_1\le r\le R_2\).
כאן \textbf{אין} סיבה לפסול אף אחד מהפתרונות הרדיאליים, כיוון שהראשית והאינסוף אינם בתחום.

\begin{resultbox}{הפתרון הכללי בטבעת}
\[
\begin{aligned}
u(r,\theta)=\;&
A_0 + B_0\ln r \\
&+ \sum_{n=1}^\infty
\Big[
(A_n r^n + B_n r^{-n})\cos(n\theta)
+ (C_n r^n + D_n r^{-n})\sin(n\theta)
\Big].
\end{aligned}
\]
\end{resultbox}

המקדמים נקבעים משני תנאי שפה: \(u(R_1,\theta)=G_1(\theta)\) ו-\(u(R_2,\theta)=G_2(\theta)\).

\subsection{בעיית גזרה זוויתית (Wedge)}
תחום: \(0\le r\le R,\quad 0\le\theta\le\alpha\).
תנאי שפה הומוגניים על הדפנות הישרות: \(u(r,0)=u(r,\alpha)=0\).

המשוואה הזוויתית עם תנאי שפה אלו גוררת ערכים עצמיים התלויים ב-\(\alpha\):
\[
\lambda_n=\frac{n\pi}{\alpha},
\qquad n=1,2,\dots
\]
הבסיס הזוויתי הוא \(\Theta_n(\theta)=\sin\!\left(\frac{n\pi}{\alpha}\theta\right)\).
הפתרון הרדיאלי הוא \(R_n(r)=r^{\lambda_n}\) (החזקה השלילית נפסלת בראשית).

\begin{resultbox}{הפתרון הכללי בגזרה}
\[
u(r,\theta)=
\sum_{n=1}^\infty
A_n
r^{\frac{n\pi}{\alpha}}
\sin\!\left(\frac{n\pi}{\alpha}\theta\right).
\]
\end{resultbox}

המקדמים \(A_n\) נקבעים על ידי הטלה (פורייה) של תנאי השפה ב-\(R\) על הבסיס \(\sin(n\pi\theta/\alpha)\).

\subsection{תובנה פיזיקלית: אפקט החוד}
כאשר \(\alpha\) קטנה (חוד חד), החזקה \(r^{\pi/\alpha}\) גדולה. אולם השדה החשמלי (הנגזרת) מתנהג כמו \(E \sim r^{\frac{\pi}{\alpha}-1}\). אם הגזרה היא \emph{המשלימה} לחומר מוליך (כלומר הזווית של המוליך היא \(2\pi-\alpha\)), אזי ליד חודים חדים השדה מתעצם מאוד.

\begin{warningbox}{בחירת פתרונות}
לעולם אין לזרוק איברים ``באוטומט''. הבחירה אילו פתרונות לפסול (מתבדרים בראשית, מתבדרים באינסוף, לוגריתמיים) תלויה אך ורק בפיזיקה של הבעיה הספציפית ובתנאי השפה והתחום המוגדרים.
\end{warningbox}
